\documentclass{article}
\usepackage{amsmath}
\usepackage{amsthm}
\usepackage{amsfonts}
\usepackage{amssymb}
\usepackage{bbold}
\usepackage{dsfont}
\usepackage{graphicx}
\usepackage{cleveref}
\usepackage{bm}
\usepackage{booktabs}
\usepackage{siunitx}
\usepackage{xcolor}
\usepackage[utf8]{inputenc}
\usepackage[english]{babel}
\newtheorem{definition}{Definition}
\newtheorem{example}{Example}
\newcommand{\exP}[1]{\mathrm{e}^{#1}}
\newcommand{\dif}{\mathrm{d}}
\renewcommand{\vec}[1]{\mathbf{#1}}
\newcommand{\pdagger}{{\phantom{\dagger}}}
\newcommand{\state}[1]{\vert #1\rangle}
\newcommand{\cstate}[1]{\langle #1\vert}
\newcommand{\tr}[1]{\mathrm{tr}\left\{#1\right\}}
\newcommand{\trs}[2]{\, \mathrm{tr}_{#1}\{#2\}}
\newcommand{\expc}[3]{\langle #1\vert #2 \vert #3\rangle}
\usepackage[doi=false,url=false,isbn=false,backend=bibtex,style=nature,natbib=true,sorting=none]{biblatex}
\bibliography{bib.bib}
\usepackage{geometry}
\geometry{
	a4paper,
	total={170mm,257mm},
	left=20mm,
	top=20mm,
}
\begin{document}
\title{Additional results and analyses: Transport through spinfull two-site Peierls model}
\author{Rudolf Smorka}
\maketitle
%\tableofcontents
\section{Model}
\subsection{Original model}
The original two-site Peierls-model is extended by having two spin states per site
\begin{equation}
	H_\mathrm{Peierls} = \sum_{i\sigma}\varepsilon_{i\sigma}^\pdagger c_{i\sigma}^\dagger c_{i\sigma}^\pdagger + (-t_0-t_1(a^\dagger + a))\sum_\sigma (c_{1\sigma}^\dagger c_{2\sigma}^\pdagger + \mathrm{h.c.}).
\end{equation}
As is assumed in Franssons work, $\varepsilon_{i\sigma}=\varepsilon$. The formulation of the leads remains as before, except the additional spin index which distinguishes the different spin-states of the leads electrons. However, the system-lead coupling now is spin-dependent
\begin{equation}
	H_\mathrm{sys-lead} = \sum_{k_L,k_R,\sigma\sigma'} (\nu_{k_L\sigma }^\pdagger d_{k_L\sigma}^\dagger c_{1\sigma'}^\pdagger + \nu_{k_R\sigma }^\pdagger d_{k_R\sigma}^\dagger  c_{2\sigma'}^\pdagger + \mathrm{h.c.}).
\end{equation}
In the wide-band limit, the coupling matrix $\Gamma$ is then
\begin{equation}
	\Gamma_{i \sigma\sigma'}^{(L,R)} = 2\pi \sum_{k_{L,R}} \vert \nu_{k_{L,R} \sigma}\vert^2 \delta(\varepsilon_{k_{L,R}\sigma}-\varepsilon_{i\sigma'})
\end{equation}
where $\varepsilon_{i\sigma'}$ is a system eigenenergy with spin $\sigma$. As before, we assume that $\nu_{k\sigma\sigma'} = \nu_k \delta_{\sigma\sigma'}$. 
\subsection{With spin-orbit coupling}
Even for the two-site system, a spin-orbit coupling term can be introduced. A minimal model has been proposed in \cite{Gutierrez_Díaz_Naaman_Cuniberti_2012}. The Peierls-Hamiltonian with this SOC then reads
\begin{equation}
	H_\mathrm{Peierls+SOC} = H_\mathrm{Peierls} + \sum_{m,n=1}^2( c_{m\uparrow}^\dagger W_{m,n}^\pdagger c_{n\downarrow}^\pdagger + c_{m\downarrow}^\dagger W^\times_{m,n} c_{n\uparrow}^\pdagger).
\end{equation}
The coupling matrix $W$ has following elements
\begin{align}
	W_{nn} &= -\alpha f(n\Delta z)\\
	W_{n,n+1} &= \frac{\alpha}{2\Delta z}\Psi(n\Delta z)\\
	W_{n,m}^\times &= -(W_{m,n})^*\\
	W_{n,n}^\times &= (W_{n,n})^*\\
	f(z) &= \partial_z \Psi(z)\\
	\Psi(z) &= \sum_{l=-1/2}^{1/2}\sum_{j=1}^2 \exP{-iQ j\Delta z} g_{l,j}(z)\\
	Q &=\frac{2\pi}{b}\\	
	g_{l,j}(z) &= \frac{1}{\sqrt{\left(1+\left[\frac{z-lb-j\Delta z}{a}\right]^2\right)^3}}.
\end{align}
Here, $\alpha=\hbar \lambda E_0$ is the spin-orbit coupling parameter, $\lambda = \frac{e\hbar}{(2mc)^2}$ the SOC strength, $E_0$ the local field of the charge distribution, $b$ is the helix pitch and $a$ the helix radius. The authors use $\alpha = \SI{5}{\milli\electronvolt \nano\meter}$.
\section{Diagonalization}
\subsection{Quantum vibration}
\subsection{Classical vibration}
\subsubsection{Peierls}
\subsubsection{Peierls+SOC}
\section{Mean-field potential}
\subsection{Peierls}
\subsection{Peierls+SOC}
\end{document}




